\documentclass[a5paper,10pt]{article}

% https://github.com/InDevRus/filippov-solutions

% Нумерация страниц
\pagenumbering{gobble}

% Настройка полей
\usepackage{geometry}
\geometry{left=1cm}
\geometry{right=1cm}
\geometry{top=1cm}
\geometry{bottom=1.5cm}

% Вставка таблицы
\usepackage{array}
\usepackage{tabularx}

% Инструменты выравнивания формул
\usepackage{amsmath}
\usepackage{mathtools}

% Диакритика в математике
\usepackage{amsfonts}

% Вычисления размеров на ходу
\usepackage{calc}

% Удобные записи производных
\usepackage[italic=true]{derivative}

% Настройки сносок
\usepackage{enotez}

% Длинные знаки векторов
\usepackage{anyfontsize}
\usepackage[b]{esvect}

% Вставка картинок
\usepackage{graphicx}

% Вставка красной строки
\usepackage{indentfirst}
\setlength{\parindent}{1em}

% Условия в командах
\usepackage{xifthen}

% Луа-скрипты
\usepackage{luacode}

% Настройка команд отдельно в математическом и текстовом режимах
\usepackage{mathcommand}

% Для повторения LaTeX кода
\usepackage{multido}

% Конфигурация отступов формул
\delimitershortfall-1sp
\usepackage{mleftright}
\mleftright

% Растягивание объектов
\usepackage{relsize}

% Обтекание текстом
\usepackage{wrapfig}
\usepackage{subcaption}
\usepackage{floatflt}

% Поддержка ввода кириллицы
\usepackage[english, russian]{babel}

% Настройка корневой позиции для компилятора
\usepackage{import}
\providecommand*{\ProjectRootPath}{..}

\import{\ProjectRootPath/preambles/macros/}{theorems}

\import{\ProjectRootPath/preambles/fonts}{font_setup}

% Знак градуса
\usepackage{siunitx}

\import{\ProjectRootPath/preambles/macros/}{operators}
\import{\ProjectRootPath/preambles/macros/}{redefinitions}

\import{\ProjectRootPath/preambles/fonts}{letters_setup}
\import{\ProjectRootPath/preambles/fonts/adjustments}{custom_superscripts}

\import{\ProjectRootPath/preambles/custom}{formatting}
\import{\ProjectRootPath/preambles/custom}{frames}
\import{\ProjectRootPath/preambles/custom}{list_setup}

% TODO: перевести команды на современный стиль объявления
\input{../preambles/plotting}

\begin{document}

Уравнение $\ddot{x}\br{t} = 3\pow{x}{2}$
не меняет вид от замены независимой переменной
$t = -u$,
так что поведение решений при $t \to +\infty$ и при $t \to -\infty$
будет одинаковым.

$\tsys{& \dot{x}\br{t} = y \\& \dot{y}\br{t} = 3\pow{x}{2}}$;
$\der{y}{x} = \dfrac {3\pow{x}{2}} {y}$.
$6\pow{x}{2} = 2yy\'$,
$\pow{y}{2} = 2\pow{x}{3} + C$.

Единственной особой точкой будет $E\br{0; 0}$.
$\der{y}{x} = \dfrac {O\br{\pow{x}{2}}} {y}$.
$A = \begin{pmatrix} 0 & 1 \\ 0 & 0 \end{pmatrix}$.
$\pow{\lambda}{2} = 0$. $\sub{\lambda}{1} = \sub{\lambda}{2} = 0$.
Таким образом, $E\br{0; 0}$ не относится к обычным особым точкам.

\begin{floatingfigure}{0.4\textwidth}
    \begin{tikzpicture}
        \pgfplotsset{
            Graph/.style = {
                thin,
                samples = 300,
                restrict y to domain = -5 : 5,
            },
            DirectionGraph/.style = {
                samples = 20,
                quiver = {scale arrows = 0.075},
                restrict y to domain = -5 : 5,
            },
            DirectionGraph1/.style = {
                DirectionGraph,
                quiver = {
                    u = {dot_x(x, y) / sqrt(dot_x(x, y) ^ 2 + dot_y(x, y) ^ 2)},
                    v = {dot_y(x, y) / sqrt(dot_x(x, y) ^ 2 + dot_y(x, y) ^ 2)}
                }
            },
            DirectionGraph2/.style = {
                DirectionGraph,
                quiver = {
                    u = {-dot_x(x, y) / sqrt(dot_x(x, y) ^ 2 + dot_y(x, y) ^ 2)},
                    v = {-dot_y(x, y) / sqrt(dot_x(x, y) ^ 2 + dot_y(x, y) ^ 2)}
                }
            },
        }
        \begin{axis}[
            trig format plots = rad,
            width = 6.25cm,
            height = 10cm,
            ylabel=$\dot{x}$,
            ymin = -1.9,
            ymax = 1.9,
            xmin = -1.2,
            xmax = 1.6, 
            xtick = {-10, ..., 10},
            ytick = {-10, ..., 10},
            minor tick num = 3,
            declare function = {
                f(\C, \x) = sqrt(2 * x ^ 3 + \C);
            }]

            \pgfplotsinvokeforeach {0, 0.1, 0.4, 0.8, 1.5} {
                \foreach \s in {-1, 1} {
                    \addplot[Graph, domain = -(#1 / 2) ^ (1 / 3) - 0.001 : -(#1 / 2) ^ (1 / 3) + 0.01]
                    {\s * f(#1, x)};
                    \addplot[Graph, domain = -(#1 / 2) ^ (1 / 3) + 0.01 : 1.9]
                    {\s * f(#1, x)};
                }
            }
            
            \pgfplotsinvokeforeach {-0.5, -1.5, -2.5} {
                \foreach \s in {-1, 1} {
                    \addplot[Graph, domain = (-#1 / 2) ^ (1 / 3) - 0.001 : (-#1 / 2) ^ (1 / 3) + 0.005]
                    {\s * f(#1, x)};
                    \addplot[Graph, domain = (-#1 / 2) ^ (1 / 3) + 0.005 : 1.9]
                    {\s * f(#1, x)};
                }
            }
        \end{axis}
    \end{tikzpicture}
\end{floatingfigure}

Следовательно, всякое решение задачи Коши 
$\ddot{x}\br{t} = 3\pow{x}{2}$,
$x\br{\sub{t}{0}} = a$,
$\dot{x}\br{\sub{t}{0}} = b$
при $t \to +\infty$ будет стремиться к $+\infty$,
кроме, случаев $\pow{b}{2} = 2\pow{a}{3}$, когда возможно стремление к началу координат.

Более точно, $\pow{b}{2} = 2\pow{a}{3}$ даёт возможность решить уравнение в квадратурах; 
сами решения \linebreak имеют форму
$x\br{t} = \br{C \pm \frac {\mathlarger{t}} {\sqrt{2}}}^{-2}$. 
Из этого видно, что решения такого вида не всегда продолжимы на $t \to +\infty$. 
При этом, если решение конкретной задачи Коши продолжимо, то оно стремится к нулю, иначе 
$x \to +\infty$ при некотором конечном $t$.

Неустойчивость нулевого решения системы \linebreak
$\tsys{& \dot{x}\br{t} = y \\& \dot{y}\br{t} = 3\pow{x}{2}}$
также можно продемонстировать, \linebreak используя функцию Ляпунова $v\br{x; y} = xy$ в области
$V = \cbr{x > 0; y > 0}$, для которой
$\der{v}{t} \Bigg| = \pow{y}{2} + 3\pow{x}{3} > 0$ внутри области $V$.

\end{document}
